% source: https://github.com/pfradin/luadraw/discussions/185#discussioncomment-17998535 (pfradin, block 1)
\documentclass[12pt]{article}
\usepackage[top=15mm,bottom=20mm,left=2cm,right=2cm]{geometry}
\usepackage[svgnames]{xcolor}
\usepackage[3d]{luadraw}
\usepackage{array}
\usepackage{esvect}
\usepackage{amsmath,amsthm}
\theoremstyle{definition}
 \newtheorem{EX}{Example}
 \newtheorem*{sol}{\textit{\textbf{Solution}}}
   \usepackage{enumitem}
  \usepackage{float}
%%%%%%%%%%%%%
% uses luacas and xintfrac
\usepackage{luacas}
\usepackage{xintfrac}

\directlua{M = luadraw.pt3d.M}% M is now global because it is used in TeX macros
\def\shortcuts{% included in each luadraw environment
local ld = luadraw
local graph3d = ld.graph3d
local pt3d = ld.pt3d
local Origin, vecI, vecJ, vecK = pt3d.Origin, pt3d.vecI, pt3d.vecJ, pt3d.vecK
local pyramid = ld.pyramid
local dproj3d = ld.dproj3d
local proj3d = ld.proj3d
local Hiddenlinestyle = ld.Hiddenlinestyle
local Hiddenlines = ld.Hiddenlines
local mShadedHidden = ld.mShadedHidden
local nearest = ld.nearest
local strReal = ld.strReal
local simplifyFrac = ld.simplifyFrac
local rad = ld.rad
local getbounds3d = ld.getbounds3d
local plane = ld.plane
}%

%% TeX  macros for distance calculs, 
%% The parameters must be of the form M(x,y,z) only, luacas must parse and convert the numbers x, y and z.
%% The formulas are explained so that the calculations can be performed by luacas.
\def\distpointpoint#1#2{%
\begin{CAS}
vars('y')
local A, B = #1, #2
local d = (B.x-A.x)^2+(B.y-A.y)^2+(B.z-A.z)^2
local a = simplify( SqrtExpression( d, 2 ) )
tex.print(a:tolatex())
\end{CAS}
}%

\def\distpointplane#1#2#3{%
\begin{CAS}
vars('y')
local A, B, n = #1, #2, #3
local d = ( (B.x-A.x)*n.x + (B.y-A.y)*n.y + (B.z-A.z)*n.z)^2 / (n.x^2+n.y^2+n.z^2)
local a = simplify( SqrtExpression( d, 2 ) )
tex.print(a:tolatex())
\end{CAS}
}%

\def\distpointline#1#2#3{%
\begin{CAS}
vars('y')
local A, B, v = #1, #2, #3
local u = M(A.x-B.x,A.y-B.y,A.z-B.z) % M must be global
local n = M(u.y*v.z-v.y*u.z, u.z*v.x-v.z*u.x, u.x*v.y-v.x*u.y)
local d = (n.x^2+n.y^2+n.z^2) / (v.x^2+v.y^2+v.z^2)
local a = simplify( SqrtExpression( d, 2 ) )
tex.print(a:tolatex())
\end{CAS}
}%

\def\distlineline#1#2#3#4{%
\begin{CAS}
vars('y')
local A, u, B, v = #1, #2, #3, #4
local n = M( u.y*v.z-v.y*u.z, u.z*v.x-v.z*u.x, u.x*v.y-v.x*u.y)
local d =  ((B.x-A.x)*n.x+(B.y-A.y)*n.y+(B.z-A.z)*n.z)^2 / (n.x^2+n.y^2+n.z^2)
local a = simplify( SqrtExpression( d, 2 ) )
tex.print(a:tolatex())
\end{CAS}
}%

\def\anglevectorvector#1#2{%
\begin{CAS}
vars('y')
local u, v = #1, #2
local d = (u.x*v.x + u.y*v.y + u.z*v.z)^2/ ((u.x^2 + u.y^2 + u.z^2)*(v.x^2 + v.y^2 + v.z^2))
local a = simplify( arccos( SqrtExpression(d,2) ) )
tex.print(a:evaluate():tolatex())
\end{CAS}
}%
% end of TeX macros
%%%%%%%%%%%

%% results formatting function
\begin{luacode}
\shortcuts
function frac(val, max)  -- converts val to a fraction " sign A / B"
    max = max or 10000 -- to stop if the process is too long
    if val == nil then return "" end
    local sg, den =  "", 1
    if val < 0 then sg = "-"; val = -val end
    local num = val
    while (den <= max) and (math.abs(num-nearest(num))>1e-8) do den = den+1; num = val*den end
    if den > max then 
        return sg..strReal(val) -- failure
    elseif den == 1 then
        return sg..nearest(num)
    else
        return sg..ld.nearest(num).."/"..den
    end
end

function trigfrac(val, max)  -- converts val (radians) to a fraction of pi " sign AxPi / B"
    max = max or 10000 -- to stop if the process is too long
    if val == nil then return "" end
    local sg, den =  "", 1
    val =val*ld.rad
    if val < 0 then sg = "-"; val = -val end
    local num = val
    while (den <= max) and (math.abs(num-nearest(num))>1e-8) do den = den+1; num = val*den end
    if den > max then 
        return false -- failure
    else
        den = den*180
        num, den = simplifyFrac(nearest(num),den)
        if num == 0 then return 0
        else
            if num == 1 then num = "" end
            return sg.."\\frac{"..num.."\\pi}{"..den.."}"
        end
    end
end

function tofrac(A) -- converts the coordinates of A
    return "M("..frac(A.x)..","..frac(A.y)..","..frac(A.z)..")"
end

function distpointpoint3d(A, B)
    return "\\distpointpoint{"..tofrac(A).."}{"..tofrac(B).."}"
end

function distpointplane(point, plane)
    local A, n = table.unpack(plane)
    return "\\distpointplane{"..tofrac(point).."}{"..tofrac(A).."}{"..tofrac(n).."}"
end

function distpointline3d(point, line)
    local A, u = table.unpack(line)
    return "\\distpointline{"..tofrac(point).."}{"..tofrac(A).."}{"..tofrac(u).."}"
end

function distlineline3d(line1, line2)
    local A, u = table.unpack(line1)
    local B, v = table.unpack(line2)
    return "\\distlineline{"..tofrac(A).."}{"..tofrac(u).."}{"..tofrac(B).."}{"..tofrac(v).."}"
end

function anglevectorvector3d(u,v, degree)
    if degree == nil then degree = true end
    local alpha = pt3d.angle(u,v)
    if degree then
        return ld.strReal(alpha*rad).."^\\circ"
    end
    local alpha = trigfrac( alpha )
    if alpha == false then 
        return "\\anglevectorvector{"..tofrac(u).."}{"..tofrac(v).."}"
    else
        return alpha
    end
end

function anglelineline3d(line1, line2, degree)
    local A, u = table.unpack(line1)
    local B, v = table.unpack(line2)
    return anglevectorvector3d(u,v, degree)
end

function anglelineplane3d(line, plane, degree)
    if degree == nil then degree = true end
    local A, u = table.unpack(line)
    local B, v = table.unpack(plane)
    local alpha = pt3d.angle(u,v)
    if degree then
        return ld.strReal(90-alpha*rad).."^\\circ"
    end
    local alpha = math.pi/2-alpha
    local alpha = trigfrac( alpha )
    if alpha == false then 
        return "\\frac{\\pi}2-\\anglevectorvector{"..tofrac(u).."}{"..tofrac(v).."}"
    else
        return alpha
    end
end

function angleplaneplane3d(plane1, plane2, degree)
    local A, u = table.unpack(plane1)
    local B, v = table.unpack(plane2)
    return anglevectorvector3d(u,v,degree)
end

function graph3d:Dcoord3d(...)
    local args = {}
    local name, dot
    for k, aux in ipairs{...} do
        if k%3 == 1 then
            name = aux
        elseif k%3 == 2 then
            dot = aux
        else
            table.append(args,{"$"..name.."(\\xintTeXsignedFrac {"..frac(dot.x).."},\\xintTeXsignedFrac {"..frac(dot.y).."},\\xintTeXsignedFrac {"..frac(dot.z).."})$", dot, aux})
        end
    end
    self:Dlabel3d(table.unpack(args))
end

function defmac(name,body)
    token.set_macro(name,body,"global")
end
\end{luacode}

\begin{document}
\begin{EX}
\begin{figure}[H]
  \centering
\begin{luadraw}{name=test}
\shortcuts
local A,B,C,D,S = M(0,0,0),M(1,0,0),M(1,1,0),M(0,1,0),M(0,0,1)
local G = (A+B+S)/3
local P = pyramid( {A,B,C,D}, S)
local K = dproj3d(A,{S,S-D})
local lineSD = {S, D-S} 
local Q = {S, pt3d.prod(S-B,S-D)}
local H = proj3d(G,Q)
local x1,x2,y1,y2,z1,z2 = getbounds3d( {A,B,C,D,S} ) -- 3d bounding box
local dl = math.max(x2-x1,y2-y1,z2-z1) + 1 -- the greatest length, plus one

local g = graph3d:new{window3d={x1,x2+1,y1,y2+1,z1,z2+1}, adjust2d=true, viewdir={20,70},size={12,12}}
Hiddenlinestyle = "dashed";Hiddenlines = true;

g:Dpolyline3d({{x1*vecI,(x1+dl)*vecI},{y1*vecJ,(y1+dl)*vecJ},{z1*vecK,(z1+dl)*vecK}}, "arrows={-Stealth[scale=1.5]}") -- axes
g:Dpoly(P, {color="orange", mode=mShadedHidden})
 g:Dlabel3d("$A$",A,{pos="S"}, "$G$",G,{},"$x$",(x1+dl)*vecI,{pos="S"},"$y$",(y1+dl)*vecJ,{pos="S"},"$z$",(z1+dl)*vecK,{pos="E"})
g:Ddots3d({A,B,C,D,S,G})
g:Dcoord3d("S",S,{pos="E"}, "C",C,{pos="S"}, "B",B,{pos="W"}, "D",D,{pos="NE"})
digits = 3
defmac("GD", distpointpoint3d(G,D))
defmac("AB", distpointpoint3d(A,B))
defmac("AD", distpointpoint3d(A,D))
defmac("AS", distpointpoint3d(A,S))
defmac("dASBD", distpointplane(A,Q) )
defmac("dASD", distpointline3d(A,lineSD) )
defmac("angleSCBD", anglelineline3d({S,C-S}, {B,D-B}) )
defmac("angleSCDSBC", angleplaneplane3d(plane(S,C,D), plane(S,B,C)) )
defmac("angleSCSAB", anglelineplane3d({S,S-C},plane(S,A,B)) )
defmac("angleACSD", anglevectorvector3d(C-A,D-S) )
defmac("angleACAB", anglevectorvector3d(C-A,B-A) )
defmac("angleSCSD", anglevectorvector3d(B-S,C-S) )
digits = 4 -- default value
g:Show()
\end{luadraw}
\caption{}
\end{figure}
Consider the pyramid $S.ABCD$ whose base $ABCD$ is a rectangle with $AB = \AB$ and $AD = \AD$. The line $SA$ is perpendicular to the plane $(ABCD)$, and $AS = \AS$. Let $G$ be the centroid of triangle $SAB$.
\begin{enumerate}[label=\arabic*)]
  \item Calculate the distance between two points $G$ and $D$.
  \item Calculate the distance from the point $A$ and the line $SD$.
  \item Calculate the distance from the point $A$ and the plane $(SBD)$.
  \item Calculate the angle between two lines $SC$ and $BD$.
  \item Calculate the angle between two plane $(SBC)$ and $(SBD)$.
  \item Calculate the angle between the line $SC$ and the plane $(SAB)$.
  \item Calculate the angle between two vectors $AC$ and  $SD$.
\end{enumerate}
\begin{sol}
\mbox{}
\begin{enumerate}[label=\arabic*)]
  \item The distance between two points $G$ and $D$ is $\mathrm d(G,D) = \GD$.
  \item The distance from the point $A$ and the line $SD$ is $\mathrm d(A, SD) = \dASD$.
  \item The distance from the point $A$ and the plane $(SBD)$ is $\mathrm d(A,(SBD)) = \dASBD$.
  \item The angle between two lines $SC$ and $BD$ is $(SC,BD) = \angleSCBD$.
  \item The angle between two plane $(SBC)$ and $(SBD)$ $((SCD),(SBC)) = \angleSCDSBC$.
\item The angle between the line $SC$ and the plane $(SAB)$ is $(SC,(SAB)) = \angleSCSAB$.
 \item   The angle between two vectors $AC$ and  $SD$ is $(\vv{AC},\vv{SD}) = \angleACSD$.
\end{enumerate}
\end{sol}
\end{EX}

\begin{EX}
\begin{figure}[H]
  \centering
\begin{luadraw}{name=test}
\shortcuts
local A,B,C,D,S = M(0,0,0),M(2,0,0),M(2,3,0),M(0,3,0),M(0,0,6)
local G = (A+B+S)/3
local P = pyramid( {A,B,C,D}, S)
local K = dproj3d(A,{S,S-D})
local lineSD = {S, D-S} 
local Q = {S, pt3d.prod(S-B,S-D)}
local H = proj3d(G,Q)
local x1,x2,y1,y2,z1,z2 = getbounds3d( {A,B,C,D,S} ) -- 3d bounding box
local dl = math.max(x2-x1,y2-y1,z2-z1) + 1 -- the greatest length, plus one

local g = graph3d:new{window3d={x1,x2+1,y1,y2+1,z1,z2+1}, adjust2d=true, viewdir={20,70},size={12,12}}
Hiddenlinestyle = "dashed";Hiddenlines = true;

g:Dpolyline3d({{x1*vecI,(x1+dl)*vecI},{y1*vecJ,(y1+dl)*vecJ},{z1*vecK,(z1+dl)*vecK}}, "arrows={-Stealth[scale=1.5]}") -- axes
g:Dpoly(P, {color="orange", mode=mShadedHidden})
 g:Dlabel3d("$A$",A,{pos="S"}, "$G$",G,{},"$x$",(x1+dl)*vecI,{pos="S"},"$y$",(y1+dl)*vecJ,{pos="S"},"$z$",(z1+dl)*vecK,{pos="E"})
g:Ddots3d({A,B,C,D,S,G})
g:Dcoord3d("S",S,{pos="E"}, "C",C,{pos="S"}, "B",B,{pos="W"}, "D",D,{pos="NE"})
digits = 3
defmac("GD", distpointpoint3d(G,D))
defmac("AB", distpointpoint3d(A,B))
defmac("AD", distpointpoint3d(A,D))
defmac("AS", distpointpoint3d(A,S))
defmac("dASBD", distpointplane(A,Q) )
defmac("dASD", distpointline3d(A,lineSD) )
defmac("angleSCBD", anglelineline3d({S,C-S}, {B,D-B}) )
defmac("angleSCDSBC", angleplaneplane3d(plane(S,C,D), plane(S,B,C)) )
defmac("angleSCSAB", anglelineplane3d({S,S-C},plane(S,A,B)) )
defmac("angleACSD", anglevectorvector3d(C-A,D-S) )
defmac("angleACAB", anglevectorvector3d(C-A,B-A) )
defmac("angleSCSD", anglevectorvector3d(B-S,C-S) )
digits = 4 -- default value
g:Show()
\end{luadraw}
\caption{}
\end{figure}
Consider the pyramid $S.ABCD$ whose base $ABCD$ is a rectangle with $AB = \AB$ and $AD = \AD$. The line $SA$ is perpendicular to the plane $(ABCD)$, and $AS = \AS$. Let $G$ be the centroid of triangle $SAB$.
\begin{enumerate}[label=\arabic*)]
  \item Calculate the distance between two points $G$ and $G$.
  \item Calculate the distance from the point $A$ and the line $SD$.
  \item Calculate the distance from the point $A$ and the plane $(SBD)$.
  \item Calculate the angle between two lines $SC$ and $BD$.
  \item Calculate the angle between two plane $(SBC)$ and $(SBD)$.
  \item Calculate the angle between the line $SC$ and the plane $(SAB)$.
  \item Calculate the angle between two vectors $AC$ and  $SD$.
\end{enumerate}
\begin{sol}
\mbox{}
\begin{enumerate}[label=\arabic*)]
  \item The distance between two points $G$ and $D$ is $\mathrm d(G,D) = \GD$.
  \item The distance from the point $A$ and the line $SD$ is $\mathrm d(A, SD) = \dASD$.
  \item The distance from the point $A$ and the plane $(SBD)$ is $\mathrm d(A,(SBD)) = \dASBD$.
  \item The angle between two lines $SC$ and $BD$ is $(SC,BD) = \angleSCBD$.
  \item The angle between two plane $(SBC)$ and $(SBD)$ $((SCD),(SBC)) = \angleSCDSBC$.
\item The angle between the line $SC$ and the plane $(SAB)$ is $(SC,(SAB)) = \angleSCSAB$.
 \item   The angle between two vectors $AC$ and  $SD$ is $(\vv{AC},\vv{SD}) = \angleACSD$.
\end{enumerate}
\end{sol}
\end{EX}
\end{document} 
